\editionPage[\pdfbookmark[0]{Défis d’approfondissement}{bm-defis-debut-parent}]{defis-debut}{Dix-huit défis}{Un défi par chapitre · Parcours facultatif}{Chercher avant de consulter les deux indications.}
\begin{blocEdition}{defis-debut}
Ces problèmes demandent davantage d’initiative. Ils utilisent les outils du livre ; les prolongements sont définis dans l’énoncé. Aucun défi n’ajoute une exigence officielle au brevet.\par\bigskip
\begin{savoirfaire}[Une recherche qui avance]Faire un schéma, tester un petit cas, comparer deux méthodes, conserver les essais et préciser ce qui reste à prouver. Ouvrir la première indication seulement après une recherche réelle ; revenir au problème avant de lire la seconde.\end{savoirfaire}
\begin{savoirfaire}[Après la solution]Identifier l’idée décisive et la refaire sur une feuille vierge. Expliquer pourquoi elle fonctionne compte autant que retrouver le résultat.\end{savoirfaire}
\end{blocEdition}
\editionPage{defi-ex-1}{D01 · Défi}{Une fraction entre deux autres}{Approfondissement · 30 à 60 minutes indicatives}
\begin{blocEdition}{defi-ex-1}
Trouver une fraction de dénominateur 60 strictement comprise entre $5/12$ et $7/15$. Combien y en a-t-il ? Puis trouver une fraction comprise entre ces deux nombres, mais de dénominateur 120.
\end{blocEdition}
\editionPage{defi-ex-2}{D02 · Défi}{Comparer sans grands calculs}{Approfondissement · 30 à 60 minutes indicatives}
\begin{blocEdition}{defi-ex-2}
Comparer $2^{30}$ et $10^9$ sans développer les deux nombres. On pourra utiliser $2^{10}=1024$. Puis dire si $2^{60}$ est plus grand que $10^{18}$.
\end{blocEdition}
\editionPage{defi-ex-3}{D03 · Défi}{Une recherche entière exhaustive}{Approfondissement · 30 à 60 minutes indicatives}
\begin{blocEdition}{defi-ex-3}
Trouver tous les entiers positifs n tels que n divise 180 et que n soit un carré parfait. Justifier que la liste est complète.
\end{blocEdition}
\editionPage{defi-ex-4}{D04 · Défi}{Un motif qui cache une identité}{Approfondissement · 30 à 60 minutes indicatives}
\begin{blocEdition}{defi-ex-4}
Pour un nombre x, comparer $(x+1)(x+2)$ et $x(x+3)$. Leur différence dépend-elle de x ? En déduire une méthode pour calculer $1001\times1002-1000\times1003$.
\end{blocEdition}
\editionPage{defi-ex-5}{D05 · Défi}{Deux solutions et un contexte}{Approfondissement · 30 à 60 minutes indicatives}
\begin{blocEdition}{defi-ex-5}
Un rectangle a pour côtés x et $x+3$ cm. Un carré a pour côté $x+1$ cm. Pour quelles valeurs positives de x leurs aires diffèrent-elles de 1 cm$^2$, sans préciser laquelle est la plus grande ?
\end{blocEdition}
\editionPage{defi-ex-6}{D06 · Défi}{Promotions dans quel ordre ?}{Approfondissement · 30 à 60 minutes indicatives}
\begin{blocEdition}{defi-ex-6}
Un article coûte P euros, avec $P>20$. Le magasin applique une réduction de 20 \% et un bon de 10 €. Comparer les deux ordres possibles et montrer que la différence de prix ne dépend pas de P.
\end{blocEdition}
\editionPage{defi-ex-7}{D07 · Défi}{Une rencontre de débits}{Approfondissement · 30 à 60 minutes indicatives}
\begin{blocEdition}{defi-ex-7}
Une cuve vide reçoit 18 L/min et perd simultanément 3 L/min par une fuite constante. Après 12 min, on coupe l’arrivée d’eau ; la fuite continue seule. Au bout de combien de temps depuis le début la cuve redevient-elle vide ?
\end{blocEdition}
\editionPage{defi-ex-8}{D08 · Défi}{Une fonction cachée derrière des essais}{Approfondissement · 30 à 60 minutes indicatives}
\begin{blocEdition}{defi-ex-8}
On sait seulement que $f(-1)=3$, $f(0)=2$, $f(1)=3$. Un élève conclut que $f(x)=x^2+2$. Montrer que cette conclusion n’est pas imposée en proposant une autre formule donnant les mêmes trois images mais une image différente pour 2.
\end{blocEdition}
\editionPage{defi-ex-9}{D09 · Défi}{Trois tarifs, un choix argumenté}{Approfondissement · 30 à 60 minutes indicatives}
\begin{blocEdition}{defi-ex-9}
Pour n séances entières, de 1 à 20, trois offres coûtent A = 8n, B = $30+5n$, C = $70+2n$, en euros. Déterminer l’offre la moins chère pour chaque nombre de séances ; indiquer les égalités utiles.
\end{blocEdition}
\editionPage{defi-ex-10}{D10 · Défi}{Faire baisser une moyenne sans baisser la médiane}{Approfondissement · 30 à 60 minutes indicatives}
\begin{blocEdition}{defi-ex-10}
Proposer une série de cinq nombres dont la moyenne vaut 10 et la médiane 8. Modifier une seule valeur pour faire baisser la moyenne à 9 tout en gardant la médiane 8. Expliquer.
\end{blocEdition}
\editionPage{defi-ex-11}{D11 · Défi}{Un jeu de somme équitable ?}{Approfondissement · 30 à 60 minutes indicatives}
\begin{blocEdition}{defi-ex-11}
Deux dés équilibrés sont lancés indépendamment. Le joueur A gagne si la somme est 6 ou 8 ; B gagne si elle est 7. Sinon on recommence sans modifier les règles. Les deux joueurs ont-ils la même chance de gagner lors d’un lancer décisif ? Comparer en comptant les couples, sans nouveau formalisme.
\end{blocEdition}
\editionPage{defi-ex-12}{D12 · Défi}{Une corde dans un cercle}{Approfondissement · 30 à 60 minutes indicatives}
\begin{blocEdition}{defi-ex-12}
Un cercle a pour centre O et rayon 13 cm. Une corde AB est à 5 cm de O. H est le pied de la perpendiculaire issue de O sur AB. On admet que H est le milieu de la corde. Calculer AB et l’aire de OAB.\label{fig-defi-12}\figureGeo{defi-corde}

\end{blocEdition}
\editionPage{defi-ex-13}{D13 · Défi}{Une droite parallèle partage une aire}{Approfondissement · 30 à 60 minutes indicatives}
\begin{blocEdition}{defi-ex-13}
Dans ABC, M appartient à [AB], N à [AC] et MN est parallèle à BC. L’aire de AMN vaut 18 cm$^2$, celle de ABC vaut 50 cm$^2$. Si AB = 15 cm et AC = 20 cm, calculer AM, AN et l’aire de MBCN.\label{fig-defi-13}\figureGeo{thales-01}

\end{blocEdition}
\editionPage{defi-ex-14}{D14 · Défi}{Deux angles de visée}{Approfondissement · 30 à 60 minutes indicatives}
\begin{blocEdition}{defi-ex-14}
Deux points d’observation A et B sont à 10 m l’un de l’autre sur une même droite horizontale allant vers le pied d’une tour. B est le plus proche. Les angles d’élévation du sommet sont de 30 degrés en A et 45 degrés en B. Les visées se font au niveau du sol dans ce modèle. Calculer la hauteur au décimètre.\label{fig-defi-14}\figureGeo{defi-visees}

\end{blocEdition}
\editionPage{defi-ex-15}{D15 · Défi}{La même aire après deux changements}{Approfondissement · 30 à 60 minutes indicatives}
\begin{blocEdition}{defi-ex-15}
Un rectangle de 6 cm sur 10 cm est agrandi par une homothétie de rapport $-2$, puis réduit par une homothétie de rapport $r>0$. Quelle valeur de r rend à l’image finale son aire initiale ? Les centres peuvent être différents.\label{fig-defi-15}\figureGeo{defi-homotheties}

\end{blocEdition}
\editionPage{defi-ex-16}{D16 · Défi}{Une section d’aire imposée}{Approfondissement · 30 à 60 minutes indicatives}
\begin{blocEdition}{defi-ex-16}
Une boule de rayon 10 cm est coupée par un plan. Le disque de section doit avoir pour aire $64\pi$ cm$^2$. À quelle distance du centre placer le plan ? Y a-t-il plusieurs positions parallèles possibles ?\label{fig-defi-16}\figureGeo{defi-section}

\end{blocEdition}
\editionPage{defi-ex-17}{D17 · Défi}{Un script de recherche et sa preuve}{Approfondissement · 30 à 60 minutes indicatives}
\begin{blocEdition}{defi-ex-17}
Un script initialise s à 0, n à 1, puis répète tant que s est inférieur à 100 : ajouter n à s ; ajouter 1 à n. Quelle valeur de n possède-t-il à la sortie ? Écrire un script par blocs ou en français structuré et justifier l’arrêt.
\end{blocEdition}
\editionPage{defi-ex-18}{D18 · Défi}{Le meilleur rectangle entier}{Approfondissement · 30 à 60 minutes indicatives}
\begin{blocEdition}{defi-ex-18}
Un rectangle a des côtés entiers positifs et un périmètre de 46 cm. Déterminer les dimensions donnant l’aire maximale et prouver qu’aucun autre rectangle entier ne fait mieux.
\end{blocEdition}
